\documentclass[11pt,a4paper]{article}

\usepackage[a4paper,margin=25mm]{geometry}
\usepackage[T1]{fontenc}
\usepackage[utf8]{inputenc}
\usepackage{lmodern}
\usepackage{microtype}

\usepackage{amsmath}
\usepackage{amssymb}
\usepackage{siunitx}
\usepackage{graphicx}
\usepackage{subcaption}
\usepackage{booktabs}
\usepackage{tabularx}
\usepackage{xcolor}
\usepackage{enumitem}
\usepackage{csquotes}
\usepackage[hidelinks]{hyperref}
\usepackage[nameinlink,noabbrev]{cleveref}

\usepackage[
  backend=biber,
  style=numeric,
  sorting=none,
  giveninits=true,
  maxbibnames=6
]{biblatex}
\addbibresource{refs.bib}

\graphicspath{{figures/}}
\sisetup{detect-all}

\title{Venue-Neutral Paper Title}
\author{
  First Author\thanks{Affiliation, email@example.com}
  \and
  Second Author\thanks{Affiliation, email@example.com}
}
\date{\today}

\begin{document}

\maketitle

\begin{abstract}
% Write a concise summary of the motivation, method, key results, and main contribution.
% Keep this venue-neutral for now; many venues will later impose a word limit or
% structured abstract format.
Abstract goes here.
\end{abstract}

% Update these once the paper scope is clearer. Some venues may require keywords
% to be removed, moved, or formatted differently.
\noindent\textbf{Keywords:} tactile sensing; electrical impedance tomography; soft robotics; sensor design

\section{Figure Plan}
\label{sec:figure-plan}

% Use this section as a working plan for the paper's figures. It can be removed,
% moved to an appendix, or converted into actual figure environments once the
% venue and final paper structure are clear.
%
% Suggested notes for each figure:
% - What question does this figure answer?
% - What data, image, schematic, or result will it show?
% - What is the intended takeaway?
% - What still needs to be generated, exported, or redesigned?

\begin{itemize}
\item \textbf{Figure 1:} Overview of the proposed sensor, the two sensing
modalities, and the principal experimental evaluations. Dhruv

\item \textbf{Figure 2:} Sensor architecture and design, including the
material structure, electrode arrangement, fabrication process, and
electrical interfacing. Aditya 

\item \textbf{Figure 3:} Experimental setup and data-collection procedure,
including contact-position control, force application, and signal
acquisition. Aditya

\item \textbf{Figure 4:} Fundamental sensor characterisation, including
signal stability, repeatability, sensitivity, hysteresis, drift, and
operating force range. Dhruv

\item \textbf{Figure 5:} Signal-processing pipelines for the direct
resistive-grid and EIT sensing modalities, including preprocessing,
feature extraction, EIT reconstruction, and localisation. Aditya

\item \textbf{Figure 6:} Spatial sensing performance using direct
resistive-grid measurements, including single-contact localisation,
spatial resolution, force dependence, and multi-contact response. Aditya

\item \textbf{Figure 7:} Spatial sensing performance using EIT
measurements, including single-contact localisation, spatial resolution,
force dependence, and multi-contact response. Aditya and Dhruv

\item \textbf{Figure 8:} Direct comparison of the two sensing modalities,
including localisation accuracy, spatial separability, measurement
superposition, multi-touch performance, and modality-specific failure
modes. Dhruv

\item \textbf{Figure 9:} Optional extension investigating distributed
contact-map estimation from the single-contact localisation dataset.

\item \textbf{Figure 10:} Optional demonstration of multi-contact or simple
contact-shape reconstruction using the estimated virtual-taxel
representation.


\end{itemize}

\section{Introduction}
\label{sec:introduction}

% Introduce the application problem, the sensing challenge, and the specific gap
% this paper addresses. End the introduction with a concise list or paragraph of
% contributions.

% Example contribution structure:
% \begin{itemize}[leftmargin=*]
%   \item A clear statement of the proposed sensor, dataset, model, analysis, or design method.
%   \item A concise description of the experimental validation.
%   \item A summary of the main finding and its relevance.
% \end{itemize}

Introduction text goes here.

\section{Related Work}
\label{sec:related-work}

% Summarise the closest prior work and position this paper against it. Keep the
% section selective: focus on work that directly motivates the design choices,
% experimental comparisons, and claimed contribution.

% Placeholder citation to confirm bibliography plumbing:
% Electrical impedance tomography is a broad imaging and sensing technique with a
% long methodological history~\cite{holder2004electrical}. Replace this with the
% specific papers that matter for the final venue and argument.

Related work text goes here.

\section{Methods}
\label{sec:methods}

% Describe the sensor design, fabrication, measurement setup, data collection
% protocol, preprocessing, and analysis pipeline. Use enough detail that another
% researcher can reproduce the work.

Methods overview goes here.

\subsection{Experimental Setup}
\label{sec:experimental-setup}

% Describe hardware, calibration, sampling rates, electrode arrangement,
% actuation or loading conditions, and environmental constraints.

Experimental setup text goes here.

\subsection{Data Processing}
\label{sec:data-processing}

% Describe filtering, baseline subtraction, feature extraction, model training,
% and validation splits. Define variables and metrics before using them.

Data processing text goes here.

\section{Results}
\label{sec:results}

% Present the main quantitative and qualitative findings. Prefer figures and
% tables that answer one question each.

Results text goes here.

% Example table scaffold:
% \begin{table}[htbp]
%   \centering
%   \caption{Main quantitative results.}
%   \label{tab:main-results}
%   \begin{tabular}{lcc}
%     \toprule
%     Method & Metric 1 & Metric 2 \\
%     \midrule
%     Baseline & -- & -- \\
%     Proposed & -- & -- \\
%     \bottomrule
%   \end{tabular}
% \end{table}

\section{Discussion}
\label{sec:discussion}

% Interpret the results rather than repeating them. Discuss why the method
% worked, where it failed, what assumptions matter, and how the findings should
% guide the next design iteration.

Discussion text goes here.

\subsection{Limitations}
\label{sec:limitations}

% State limitations plainly. Include dataset size, hardware constraints,
% generalisation limits, sources of uncertainty, and any venue-relevant ethical or
% reproducibility considerations.

Limitations text goes here.

\section{Conclusion}
\label{sec:conclusion}

% Close with the main contribution, the strongest result, and the most important
% next step. Keep the conclusion short unless the target venue expects a broader
% final section.

Conclusion text goes here.


\printbibliography

\end{document}
